\documentclass[twocolumn]{article}

\usepackage[utf8]{inputenc}
\usepackage{amsmath,amssymb}
\usepackage{graphicx}
\usepackage{booktabs}
\usepackage{hyperref}
\usepackage{natbib}
\usepackage[margin=1in]{geometry}
\usepackage{caption}
\usepackage{subcaption}
\usepackage{algorithm}
\usepackage{algpseudocode}

\title{Multi-Modal Deep Learning for Spacecraft Orbit Prediction:\\
  Incorporating Solar Wind Perturbations via Cross-Attention Fusion}

\author{
  Ted Rubin \\
  % \textit{Affiliation} \\
  % \texttt{email@example.com}
}

\date{\today}

\begin{document}

\maketitle

% ============================================================
\begin{abstract}
Accurate spacecraft orbit prediction is critical for collision avoidance,
mission planning, and space situational awareness. Traditional propagators
(SGP4, numerical integrators) rely on simplified atmospheric and
gravitational models that degrade during geomagnetic storms when the
thermosphere expands unpredictably. We present a multi-modal deep learning
approach that combines historical spacecraft position time series with
real-time solar wind measurements to improve orbit prediction accuracy
during space weather events. Using three years (2023--2025) of NASA SSC
position data for spacecraft in LEO, L1, and highly elliptical orbits,
combined with OMNI solar wind parameters, we demonstrate that a
cross-attention fusion architecture outperforms both SGP4 baselines and
single-modality LSTM/Transformer models, particularly during periods of
elevated geomagnetic activity (Kp $> 5$). Our results suggest that solar
wind information provides a leading indicator of orbit perturbations with
an optimal lag of [TBD] hours, enabling proactive trajectory correction.

\bigskip
\noindent\textbf{Keywords:} orbit prediction, deep learning, solar wind,
space weather, LSTM, Transformer, multi-modal learning, spacecraft tracking
\end{abstract}

% ============================================================
\section{Introduction}
\label{sec:intro}

Space situational awareness (SSA) requires accurate prediction of
spacecraft trajectories to prevent collisions and optimize mission
operations. The standard approach uses the Simplified General Perturbations
model (SGP4)~\citep{vallado2006revisiting}, which propagates orbits using
Two-Line Element (TLE) sets. While SGP4 performs well for short-term
predictions under quiet conditions, its accuracy degrades significantly
during geomagnetic storms when the thermosphere heats and expands,
increasing atmospheric drag on low Earth orbit (LEO) satellites.

Recent advances in deep learning for time-series forecasting have opened
new possibilities for orbit prediction. Unlike physics-based propagators,
neural networks can learn complex, nonlinear relationships between
environmental conditions and trajectory perturbations directly from data.

\subsection{Contributions}

This paper makes the following contributions:

\begin{enumerate}
  \item We demonstrate that LSTM and Transformer architectures can
    outperform SGP4 for short-horizon (6--24h) orbit prediction when
    trained on sufficient historical data.
  \item We introduce a cross-attention multi-modal architecture that fuses
    spacecraft position sequences with solar wind measurements, showing
    [TBD]\% error reduction during geomagnetic storms.
  \item We provide a systematic analysis of the lag relationship between
    solar wind conditions at L1 and orbit perturbations for different
    orbit types (LEO, HEO, L1).
  \item We release an open-source pipeline for fetching, processing, and
    modeling spacecraft trajectory data from NASA SSC.
\end{enumerate}

% ============================================================
\section{Related Work}
\label{sec:related}

\subsection{Physics-Based Orbit Prediction}

% SGP4, numerical propagators, atmospheric density models
% Limitations during geomagnetic storms
% [TO BE FILLED WITH LITERATURE REVIEW]

\subsection{Machine Learning for Orbit Prediction}

% Recent ML approaches to orbit determination
% LSTM for trajectory forecasting
% Transformers for time-series
% [TO BE FILLED]

\subsection{Solar Wind and Space Weather}

% Solar wind measurements at L1
% Geomagnetic indices (Kp, Dst)
% Thermosphere response to geomagnetic storms
% [TO BE FILLED]

% ============================================================
\section{Data}
\label{sec:data}

\subsection{Spacecraft Position Data}

We use the NASA Satellite Situation Center (SSC) REST
API\footnote{\url{https://sscweb.gsfc.nasa.gov/WebServices/REST/}} to
obtain spacecraft positions in the Geocentric Solar Ecliptic (GSE)
coordinate system at 1-minute resolution. We select three spacecraft
representing distinct orbital regimes:

\begin{itemize}
  \item \textbf{ISS} (International Space Station): Low Earth Orbit at
    $\sim$408~km altitude, subject to significant atmospheric drag.
  \item \textbf{DSCOVR} (Deep Space Climate Observatory): Lissajous orbit
    around the Sun-Earth L1 Lagrange point ($\sim$1.5 million km from Earth).
  \item \textbf{MMS-1} (Magnetospheric Multiscale Mission): Highly
    elliptical orbit passing through Earth's magnetosphere.
\end{itemize}

Our dataset spans January 2023 through [current date] 2025, comprising
approximately [TBD] million position records.

\subsection{Solar Wind Data}

Solar wind parameters are obtained from the OMNI database via NASA's
Coordinated Data Analysis Web (CDAWeb) service. We use the high-resolution
(1-minute) OMNI dataset, which provides:

\begin{itemize}
  \item Interplanetary Magnetic Field (IMF) components: $B_x$, $B_y$,
    $B_z$ in GSE coordinates
  \item Solar wind bulk flow speed
  \item Proton number density
  \item Geomagnetic indices: Kp, Dst
\end{itemize}

The OMNI data is time-shifted to account for solar wind propagation from
the measurement point (typically L1) to Earth's bow shock.

\subsection{Data Preprocessing}

% Velocity derivation via finite differences
% Normalization (per-spacecraft, zero-mean unit-variance)
% Gap handling and segment identification
% Sliding window creation
% Temporal train/val/test split

% ============================================================
\section{Methods}
\label{sec:methods}

\subsection{Problem Formulation}

Given a sequence of $T$ past spacecraft states
$\mathbf{X} = [\mathbf{x}_1, \ldots, \mathbf{x}_T]$ where each
$\mathbf{x}_t \in \mathbb{R}^6$ contains position and velocity components
$(x, y, z, \dot{x}, \dot{y}, \dot{z})$ in GSE coordinates, we predict
the next $H$ positions $\hat{\mathbf{Y}} = [\hat{\mathbf{y}}_{T+1},
\ldots, \hat{\mathbf{y}}_{T+H}]$ where $\hat{\mathbf{y}}_t \in
\mathbb{R}^3$ is the predicted position $(x, y, z)$.

For the multi-modal setting, we additionally condition on a concurrent
solar wind sequence $\mathbf{S} = [\mathbf{s}_1, \ldots, \mathbf{s}_T]$
where $\mathbf{s}_t \in \mathbb{R}^7$ contains IMF components, flow
speed, proton density, Kp, and Dst.

\subsection{Baseline: SGP4 / Kepler Propagation}

% Two-body Kepler propagation with Velocity Verlet integration
% SGP4 propagation from initial state

\subsection{LSTM Architecture}

We employ a bidirectional LSTM encoder with a unidirectional decoder:

\begin{equation}
  \mathbf{h}_t^{\text{enc}} = \text{BiLSTM}(\mathbf{x}_t, \mathbf{h}_{t-1}^{\text{enc}})
\end{equation}

The encoder's final hidden state is projected through a bridge layer and
used to initialize a decoder LSTM that autoregressively generates
predictions. We also evaluate a direct (non-autoregressive) variant that
produces all $H$ outputs simultaneously via a global pooling and MLP head.

\subsection{Transformer Architecture}

We use an encoder-decoder Transformer with sinusoidal positional
encodings. The encoder processes the input sequence via multi-head
self-attention, and the decoder uses learnable query tokens with
cross-attention to the encoder memory:

\begin{equation}
  \text{Attention}(Q, K, V) = \text{softmax}\left(\frac{QK^\top}{\sqrt{d_k}}\right) V
\end{equation}

\subsection{Multi-Modal Cross-Attention Fusion}

For incorporating solar wind data, we propose a dual-encoder architecture
with cross-modal attention:

\begin{enumerate}
  \item \textbf{Orbit Encoder}: Bidirectional LSTM processes spacecraft
    position/velocity sequences.
  \item \textbf{Solar Wind Encoder}: Separate bidirectional LSTM processes
    solar wind parameter sequences.
  \item \textbf{Cross-Attention}: The orbit representation attends to
    solar wind features via multi-head cross-attention, learning which
    solar wind conditions most affect the trajectory.
  \item \textbf{Fusion}: Attended orbit features and pooled solar wind
    features are concatenated and passed through an MLP prediction head.
\end{enumerate}

The model predicts position \emph{residuals} (deviations from the nominal
orbit), focusing learning capacity on the perturbation signal.

% ============================================================
\section{Experiments}
\label{sec:experiments}

\subsection{Experimental Setup}

% Training details: AdamW, cosine LR, early stopping
% Hardware: [TBD]
% Hyperparameters table

\subsection{Evaluation Metrics}

We evaluate using Mean Absolute Error (MAE) and Root Mean Square Error
(RMSE) of the 3D Euclidean position error in kilometers:

\begin{equation}
  \text{MAE} = \frac{1}{NH} \sum_{i=1}^{N} \sum_{t=1}^{H}
    \|\hat{\mathbf{y}}_{i,t} - \mathbf{y}_{i,t}\|_2
\end{equation}

We report metrics at multiple prediction horizons (1h, 6h, 24h) and
separately for quiet (Kp $\leq 3$) and storm (Kp $> 5$) conditions.

% ============================================================
\section{Results}
\label{sec:results}

% [TO BE FILLED WITH EXPERIMENTAL RESULTS]

\subsection{Orbit Prediction Performance}

% Table: Model comparison (SGP4 vs LSTM vs Transformer)
% Figure: Error vs prediction horizon
% Per-spacecraft analysis

\begin{table}[h]
\centering
\caption{Orbit prediction results (ISS, 6-hour horizon)}
\label{tab:results}
\begin{tabular}{lcccc}
\toprule
Model & MAE (km) & RMSE (km) & MAE$_{1h}$ & MAE$_{6h}$ \\
\midrule
Kepler        & -- & -- & -- & -- \\
SGP4          & -- & -- & -- & -- \\
LSTM          & -- & -- & -- & -- \\
Transformer   & -- & -- & -- & -- \\
Multi-Modal   & -- & -- & -- & -- \\
\bottomrule
\end{tabular}
\end{table}

\subsection{Solar Wind Impact Analysis}

% Cross-correlation analysis
% Optimal lag determination
% Multi-modal vs orbit-only comparison
% Storm event case studies

\subsection{Ablation Studies}

% Which solar wind parameters matter most
% Effect of input window length
% Attention visualization

% ============================================================
\section{Discussion}
\label{sec:discussion}

% Interpretation of results
% Why multi-modal helps during storms
% Limitations
% Comparison to operational systems

% ============================================================
\section{Conclusion}
\label{sec:conclusion}

% Summary of findings
% Implications for operational orbit prediction
% Future work: more spacecraft, real-time deployment, ensemble methods

% ============================================================
\section*{Data Availability}

All spacecraft position data is publicly available through the NASA
Satellite Situation Center (SSC) REST API. Solar wind data is available
through NASA's CDAWeb OMNI database. Our code and trained models are
available at \url{https://github.com/tedrubin80/OrbitResdearch}.

\bibliographystyle{plainnat}
\bibliography{references}

\end{document}
